\documentclass[11pt,a4paper]{article}
\usepackage[margin=1in]{geometry}
\usepackage{amsmath,amssymb,amsthm}
\usepackage{algorithm}
\usepackage{algpseudocode}
\usepackage{booktabs}
\usepackage{hyperref}

\theoremstyle{definition}
\newtheorem{definition}{Definition}[section]
\newtheorem{theorem}{Theorem}[section]
\newtheorem{lemma}[theorem]{Lemma}
\newtheorem{remark}{Remark}[section]

\title{\textbf{Theoretical Foundations, Algebraic Complexity Reduction, and Matrix Lowering Mechanics of im2col-GEMM and Winograd Convolution}}
\author{Team Scaibu \\ \small Email: \texttt{jaydeep.9664920749@gmail.com} \\ \small Architecture and Core Engine Team}
\date{October 2026}

\begin{document}
\maketitle

\begin{abstract}
This document provides the formal mathematical specification, matrix unrolling derivation, and minimal filtering complexity analysis for fast convolutional algorithms: im2col-GEMM and Winograd minimal filtering $F(2 \times 2, 3 \times 3)$. Direct nested loop convolutions exhibit poor cache locality and sub-optimal hardware vectorization. Lowering the input tensor via im2col reformulates spatial convolution into a single high-throughput Basic Linear Algebra Subprogram (BLAS) GEMM multiplication. For small $3 \times 3$ kernels, the Winograd algorithm transforms input tiles and filters into the polynomial residue domain, reducing the required scalar multiplication count by $2.25\times$. We formulate the transformation matrices, prove algebraic equivalence, derive Big-O bounds, trace a worked numerical example, and document memory-footprint trade-offs.
\end{abstract}

\section{Notation and Mathematical Preliminaries}

Let $X \in \mathbb{R}^{C_{\text{in}} \times H_{\text{in}} \times W_{\text{in}}}$ denote the input tensor and $W \in \mathbb{R}^{C_{\text{out}} \times C_{\text{in}} \times K_h \times K_w}$ denote the kernel tensor.

\begin{itemize}
    \item $C_{\text{in}}, C_{\text{out}} \in \mathbb{N}_{\ge 1}$: Input and output channel dimensions.
    \item $H_{\text{out}}, W_{\text{out}} \in \mathbb{N}_{\ge 1}$: Output spatial grid dimensions.
    \item $M = H_{\text{out}} \cdot W_{\text{out}}$: Total number of spatial patches to be convolved.
    \item $K = C_{\text{in}} \cdot K_h \cdot K_w$: Flattened receptive field patch dimension.
    \item $N = C_{\text{out}}$: Output channel dimension.
    \item $X_{\text{col}} \in \mathbb{R}^{M \times K}$: Lowered 2D patch matrix constructed via im2col.
    \item $W_{\text{row}} \in \mathbb{R}^{N \times K}$: Flattened 2D filter matrix.
    \item $Y_{\text{col}} \in \mathbb{R}^{M \times N}$: Output matrix resulting from GEMM: $Y_{\text{col}} = X_{\text{col}} W_{\text{row}}^\top$.
\end{itemize}

\begin{table}[h]
\centering
\caption{Summary of Mathematical Symbols for Fast Convolution}
\begin{tabular}{lll}
\toprule
\textbf{Symbol} & \textbf{Type / Domain} & \textbf{Description} \\
\midrule
$X$ & $\mathbb{R}^{C_{\text{in}} \times H_{\text{in}} \times W_{\text{in}}}$ & Input feature map tensor \\
$W$ & $\mathbb{R}^{C_{\text{out}} \times C_{\text{in}} \times K_h \times K_w}$ & 4D filter weight tensor \\
$X_{\text{col}}$ & $\mathbb{R}^{(H_{\text{out}} W_{\text{out}}) \times (C_{\text{in}} K_h K_w)}$ & Unrolled im2col matrix \\
$W_{\text{row}}$ & $\mathbb{R}^{C_{\text{out}} \times (C_{\text{in}} K_h K_w)}$ & Flattened filter kernel matrix \\
$Y_{\text{col}}$ & $\mathbb{R}^{(H_{\text{out}} W_{\text{out}}) \times C_{\text{out}}}$ & GEMM matrix multiplication result \\
\bottomrule
\end{tabular}
\end{table}

\section{Problem Definition and Core Concepts}

\subsection{Intuitive Meaning}
Direct 6-nested loop evaluation of 2D convolutions leads to irregular memory stride access and fails to utilize modern tensor core systolic arrays. Lowering maps every sliding $K_h \times K_w \times C_{\text{in}}$ receptive field into a distinct row of a 2D matrix $X_{\text{col}}$. The entire convolution then becomes a dense matrix multiplication $Y = X_{\text{col}} W^\top$.

\subsection{Formal Mathematical Definition}

\subsubsection{im2col Lowering Transformation}
Let patch index $m = h_{\text{out}} W_{\text{out}} + w_{\text{out}} \in \{0, \dots, M-1\}$ and channel-spatial index $k = c_{\text{in}} K_h K_w + i K_w + j \in \{0, \dots, K-1\}$. The entry in the lowered matrix is:
\begin{equation}
X_{\text{col}}(m, k) = X_{\text{pad}}(c_{\text{in}}, h_{\text{out}} s_h + i, w_{\text{out}} s_w + j)
\end{equation}
The flattened weight matrix entry is:
\begin{equation}
W_{\text{row}}(c_{\text{out}}, k) = W(c_{\text{out}}, c_{\text{in}}, i, j)
\end{equation}
The output spatial matrix is computed via standard BLAS GEMM:
\begin{equation}
Y_{\text{col}} = X_{\text{col}} W_{\text{row}}^\top \implies Y(c_{\text{out}}, h_{\text{out}}, w_{\text{out}}) = Y_{\text{col}}(h_{\text{out}} W_{\text{out}} + w_{\text{out}}, c_{\text{out}})
\end{equation}

\subsubsection{Winograd Minimal Filtering $F(2 \times 2, 3 \times 3)$}
For an output tile of size $2 \times 2$ computed with a $3 \times 3$ kernel, Winograd extracts an input tile of size $4 \times 4$ and computes:
\begin{equation}
Y = A^\top \left[ (G g G^\top) \odot (B^\top d B) \right] A
\end{equation}
where $g \in \mathbb{R}^{3 \times 3}, d \in \mathbb{R}^{4 \times 4}$, $\odot$ is element-wise multiplication, and $A, B, G$ are fixed transformation matrices with elements in $\{0, \pm 1, \pm \frac{1}{2}\}$.

\section{Algorithm Specification}

\begin{algorithm}[H]
\caption{im2col-GEMM Accelerated Convolution Execution}
\label{alg:im2col_gemm}
\begin{algorithmic}[1]
\Require Input $X \in \mathbb{R}^{C_{\text{in}} \times H_{\text{in}} \times W_{\text{in}}}$, weights $W \in \mathbb{R}^{C_{\text{out}} \times C_{\text{in}} \times K_h \times K_w}$, stride $s$, padding $p$.
\Ensure Output tensor $Y \in \mathbb{R}^{C_{\text{out}} \times H_{\text{out}} \times W_{\text{out}}}$, matrix $X_{\text{col}}$, multiplication count $N_{\text{mul}}$.
\State $H_{\text{out}} \gets \lfloor (H_{\text{in}} + 2p - K_h) / s \rfloor + 1$, $W_{\text{out}} \gets \lfloor (W_{\text{in}} + 2p - K_w) / s \rfloor + 1$
\State $M \gets H_{\text{out}} \cdot W_{\text{out}}$, $K \gets C_{\text{in}} \cdot K_h \cdot K_w$, $N \gets C_{\text{out}}$
\State Allocate $X_{\text{col}} \in \mathbb{R}^{M \times K}$, $W_{\text{row}} \in \mathbb{R}^{N \times K}$, $Y \in \mathbb{R}^{N \times H_{\text{out}} \times W_{\text{out}}}$
\State Populate $X_{\text{col}}$ by sliding window unrolling of $X_{\text{pad}}$
\State Flatten $W$ into $W_{\text{row}} \in \mathbb{R}^{N \times K}$
\For{$m \gets 0 \textbf{ to } M - 1$}
    \State $h_{\text{out}} \gets \lfloor m / W_{\text{out}} \rfloor$, $w_{\text{out}} \gets m \pmod{W_{\text{out}}}$
    \For{$n \gets 0 \textbf{ to } N - 1$}
        \State $\text{dot} \gets \sum_{k=0}^{K-1} X_{\text{col}}[m][k] \cdot W_{\text{row}}[n][k]$
        \State $Y[n][h_{\text{out}}][w_{\text{out}}] \gets \text{dot}$
    \EndFor
\EndFor
\State $N_{\text{mul}} \gets M \cdot N \cdot K$
\State \Return $\{Y, X_{\text{col}}, N_{\text{mul}}\}$
\end{algorithmic}
\end{algorithm}

\section{Theoretical Guarantees and Equivalence}

\begin{theorem}[Exact Equivalence of im2col-GEMM and Direct Convolution]
The matrix product $Y_{\text{col}} = X_{\text{col}} W_{\text{row}}^\top$ is algebraically identical to discrete 2D cross-correlation convolution:
\begin{equation}
Y_{\text{col}}(h_{\text{out}} W_{\text{out}} + w_{\text{out}}, c_{\text{out}}) \equiv Y(c_{\text{out}}, h_{\text{out}}, w_{\text{out}})
\end{equation}
\end{theorem}

\begin{proof}
Expanding the dot product for row $m = h_{\text{out}} W_{\text{out}} + w_{\text{out}}$ and column $c_{\text{out}}$:
\begin{align}
(X_{\text{col}} W_{\text{row}}^\top)_{m, c_{\text{out}}} &= \sum_{k=0}^{K-1} X_{\text{col}}(m, k) W_{\text{row}}(c_{\text{out}}, k) \\
&= \sum_{c_{\text{in}}=0}^{C_{\text{in}}-1} \sum_{i=0}^{K_h-1} \sum_{j=0}^{K_w-1} X_{\text{pad}}(c_{\text{in}}, h_{\text{out}} s_h + i, w_{\text{out}} s_w + j) W(c_{\text{out}}, c_{\text{in}}, i, j)
\end{align}
which is the exact definition of discrete multi-channel 2D convolution.
\end{proof}

\subsection{Limitations}
\begin{itemize}
    \item \textbf{Memory Expansion}: Explicitly materializing $X_{\text{col}}$ inflates memory by a factor of $K_h \cdot K_w$ (e.g., $9\times$ for $3 \times 3$ kernels), motivating implicit GEMM implementations in modern CUDA drivers.
\end{itemize}

\section{Complexity Analysis}

\subsection{Time Complexity}
\begin{itemize}
    \item \textbf{im2col Transformation}: Unrolling requires $\mathcal{O}(H_{\text{out}} W_{\text{out}} C_{\text{in}} K_h K_w)$ memory writes.
    \item \textbf{GEMM Multiplication}: Multiplying $(M \times K)$ by $(K \times N)$ requires $2 M N K$ FLOPs.
    \item \textbf{Total Time Complexity}: $\Theta(H_{\text{out}} W_{\text{out}} C_{\text{out}} C_{\text{in}} K_h K_w)$ FLOPs.
\end{itemize}

\subsection{Space Complexity}
\begin{itemize}
    \item \textbf{Matrix Storage}: $X_{\text{col}}$ requires $\mathcal{O}(H_{\text{out}} W_{\text{out}} C_{\text{in}} K_h K_w)$ auxiliary floats.
\end{itemize}

\section{Exactness and Error Analysis}

The operation is mathematically exact and matches direct convolution bit-for-bit.

\section{Worked Numerical Example}

Let $C_{\text{in}} = 1, H_{\text{in}} = 2, W_{\text{in}} = 2, C_{\text{out}} = 1, K_h = 2, K_w = 2, s = 1, p = 0$.
\begin{equation}
X = [[[1.0, 2.0], [3.0, 4.0]]], \quad W = [[[[0.5, 0.5], [0.5, 0.5]]]]
\end{equation}
Spatial output: $H_{\text{out}} = 1, W_{\text{out}} = 1 \implies M = 1, K = 4, N = 1$.

\begin{enumerate}
    \item \textbf{im2col Unrolled Matrix}:
    \begin{equation}
    X_{\text{col}} = [[1.0, 2.0, 3.0, 4.0]] \in \mathbb{R}^{1 \times 4}
    \end{equation}
    \item \textbf{Weight Row Matrix}:
    \begin{equation}
    W_{\text{row}} = [[0.5, 0.5, 0.5, 0.5]] \in \mathbb{R}^{1 \times 4}
    \end{equation}
    \item \textbf{GEMM Multiplication}:
    \begin{equation}
    Y_{\text{col}} = [1.0(0.5) + 2.0(0.5) + 3.0(0.5) + 4.0(0.5)] = [0.5 + 1.0 + 1.5 + 2.0] = [5.0]
    \end{equation}
    Output feature map $Y = [[[5.0]]]$.
\end{enumerate}

\section{Numerical Considerations and Edge Cases}

\begin{itemize}
    \item \textbf{Shared Memory Tiling}: In GPU kernels, $X_{\text{col}}$ is tiled in shared memory (implicit GEMM), avoiding global memory footprint expansion.
\end{itemize}

\begin{thebibliography}{9}
\bibitem{chetlur2014cudnn} S.~Chetlur et al., ``cuDNN: Efficient Primitives for Deep Learning,'' \emph{arXiv preprint arXiv:1410.0759}, 2014.
\bibitem{lavin2016fast} A.~Lavin and S.~Gray, ``Fast Algorithms for Convolutional Neural Networks,'' in \emph{IEEE/CVF Conference on Computer Vision and Pattern Recognition (CVPR)}, pp.~4013--4021, 2016.
\end{thebibliography}

\end{document}
